%!TEX root =../mapp-challenge-18-game-book.tex
% ^ leave for LaTeXTools build functionality

\phChapterWorksheet{An Unending Enigma}{Main Puzzle 4}

\newcommand{\mappDoppl}{\includegraphics[height=0.12in]{assets/doppl.pdf}}
\newcommand{\mappDopplOmega}{\includegraphics[height=0.12in]{assets/doppl-omega.pdf}}
\newcommand{\mappDopplOmegaSquared}{\includegraphics[height=0.12in]{assets/doppl-omega2.pdf}}

As you make your way down Interstate \(e\), you find a group of
\textbf{Doppl} \mappMobimon{} in a patch of tall grass.
At first the size of the group is \textbf{finite}, so you can count them
in the usual way.

\begin{itemize}
\item \(0=\) no Doppl

\item \(1=\) \mappDoppl

\item \(2=\) \mappDoppl\mappDoppl

\item \(3=\) \mappDoppl\mappDoppl\mappDoppl

\item \(4=\) \mappDoppl\mappDoppl\mappDoppl\mappDoppl
\end{itemize}

However, before you know it, the size of the group grows to
be literally \textbf{infinite}! Your \mappMobidex{} (an electronic
guide used by most trainers)
informs you that the lowercase Greek letter \(\omega\)
("omega") is used
to represent the shortest possible chain of infinitely-many Doppl.

\begin{itemize}
  \item \(\omega=\)
  \mappDoppl\mappDoppl\mappDoppl\mappDoppl\mappDoppl\mappDoppl\mappDoppl
  \mappDoppl\mappDoppl\mappDoppl\mappDoppl\mappDoppl\mappDoppl\mappDoppl
  \mappDoppl\mappDoppl\mappDoppl\mappDoppl\mappDoppl\mappDoppl\mappDoppl
  \mappDoppl\mappDoppl\mappDoppl\mappDoppl\mappDoppl\mappDoppl\mappDoppl
  \mappDoppl\mappDoppl\mappDoppl\mappDoppl\mappDoppl\mappDoppl\mappDoppl
  \mappDoppl\mappDoppl\mappDoppl\mappDoppl\mappDoppl\mappDoppl\mappDoppl
  \mappDoppl\mappDoppl\mappDoppl\mappDoppl\mappDoppl\mappDoppl\mappDoppl
\end{itemize}

As soon as they appeared, most of the Doppl \textbf{shrink} down to fill the
available space. Now, there's still an \(\omega\)-length chain of Doppl, but
they look a little more like this.

\begin{itemize}
\item \(\omega=\) \mappDopplOmega
\end{itemize}

Your \mappMobidex{} points out that while \(\omega\) is the first
\textbf{infinite number}, but it's not the only one!
You see, bizzare as it sounds, there's always room for another
Doppl to join the party.

\begin{itemize}
\item\(\omega+1=\) \mappDopplOmega\mappDoppl

\item\(\omega+2=\) \mappDopplOmega\mappDoppl\mappDoppl

\item\(\omega+3=\) \mappDopplOmega\mappDoppl\mappDoppl\mappDoppl
\end{itemize}

Suddenly, it gets really weird as you notice more infinite-length families
of Doppl in the surrounding area. Your \mappMobidex{} is able to count
a few of these groups.

\begin{itemize}
\item \(\omega+\omega=\omega\cdot2=\) \mappDopplOmega\mappDopplOmega

\item \(\omega+\omega+\omega+5=\omega\cdot 3 + 5=\) \mappDopplOmega\mappDopplOmega\mappDopplOmega\mappDoppl\mappDoppl\mappDoppl\mappDoppl\mappDoppl

\item \(\omega+\omega+\omega+\cdots=\omega\cdot\omega=\omega^2=\)
\mappDopplOmega\mappDopplOmega\mappDopplOmega\mappDopplOmega\mappDopplOmega
\mappDopplOmega\mappDopplOmega\mappDopplOmega\mappDopplOmega\mappDopplOmega
\mappDopplOmega\mappDopplOmega\mappDopplOmega\mappDopplOmega\mappDopplOmega
\mappDopplOmega\mappDopplOmega\mappDopplOmega\mappDopplOmega\mappDopplOmega
\mappDopplOmega\mappDopplOmega\mappDopplOmega\mappDopplOmega\mappDopplOmega
\mappDopplOmega\mappDopplOmega\mappDopplOmega\mappDopplOmega\mappDopplOmega
\mappDopplOmega\mappDopplOmega\mappDopplOmega\mappDopplOmega\mappDopplOmega
\mappDopplOmega\mappDopplOmega\mappDopplOmega\mappDopplOmega\mappDopplOmega

= \mappDopplOmegaSquared
\end{itemize}

\newpage

Your \mappMobidex{} points out that the rules of
\textbf{arithmetic} behave differently when the numbers involved aren't all
finite. (You suddenly remember your math teacher pleading with you to not
use \(\infty\) like it's a real number, and now realize why!)

When chains of Doppl are added
together, a finite chain of Doppl will be absorbed into an infinite chain
on its right, but not its left.

\begin{itemize}
\item \(3+4=\) \{\mappDoppl\mappDoppl\mappDoppl\}%
\{\mappDoppl\mappDoppl\mappDoppl\mappDoppl\} =
\mappDoppl\mappDoppl\mappDoppl\mappDoppl\mappDoppl\mappDoppl\mappDoppl{}
\(=7\)

\item \(7+\omega=\)
\{\mappDoppl\mappDoppl\mappDoppl\mappDoppl\mappDoppl\mappDoppl\mappDoppl\}%
\{\mappDopplOmega\} =
\mappDoppl\mappDoppl\mappDoppl\mappDoppl\mappDoppl\mappDoppl\mappDoppl%
\mappDopplOmega =
\mappDopplOmega{} =
\(\omega\)

\item \((\omega\cdot4+3)+(\omega\cdot2+5)=\)
\{\mappDopplOmega\mappDopplOmega\mappDopplOmega\mappDopplOmega\mappDoppl
\mappDoppl\mappDoppl\}%
\{\mappDopplOmega\mappDopplOmega\mappDoppl\mappDoppl\mappDoppl\mappDoppl
\mappDoppl\}

\hspace{1.55in} \(=\)
\{\mappDopplOmega\}\{\mappDopplOmega\}\{\mappDopplOmega\}\{\mappDopplOmega\}%
\{\mappDoppl\mappDoppl\mappDoppl\mappDopplOmega\}%
\{\mappDopplOmega\}\mappDoppl\mappDoppl\mappDoppl\mappDoppl\mappDoppl{}

\hspace{1.55in} \(=\omega\cdot6+5\)
\end{itemize}

The \mappMobidex{} informs you that
Doppl chains can also be multiplied. When this occurs, each Doppl in the second
factor splits into a copy of the Doppl chain given by the first factor.

\begin{itemize}
\item \(3\cdot 4=\)
\{\mappDoppl\mappDoppl\mappDoppl\}\{\mappDoppl\mappDoppl\mappDoppl\}\{\mappDoppl\mappDoppl\mappDoppl\}\{\mappDoppl\mappDoppl\mappDoppl\} \(=12\)

\item \((\omega+1)\cdot 2=\)
\{\mappDopplOmega\mappDoppl\}\{\mappDopplOmega\mappDoppl\} =
\{\mappDopplOmega\}\{\mappDoppl\mappDopplOmega\}\mappDoppl{} \(=\omega\cdot 2+1\)

\item \(2\cdot(\omega+1)=\)
\{\mappDoppl\mappDoppl\}%
\resizebox{0.4in}{!}{\{\mappDoppl\mappDoppl\}}%
\resizebox{0.3in}{!}{\{\mappDoppl\mappDoppl\}}%
\resizebox{0.2in}{!}{\{\mappDoppl\mappDoppl\}}%
\resizebox{0.1in}{!}{\{\mappDoppl\mappDoppl\}}%
\{\mappDoppl\mappDoppl\} =
\mappDopplOmega\mappDoppl\mappDoppl \(=\omega+2\)

\item \((\omega+1)\cdot \omega=\)
\{\mappDopplOmega\mappDoppl\}%
\resizebox{0.6in}{!}{\{\mappDopplOmega\mappDoppl\}}%
\resizebox{0.5in}{!}{\{\mappDopplOmega\mappDoppl\}}%
\resizebox{0.4in}{!}{\{\mappDopplOmega\mappDoppl\}}%
\resizebox{0.3in}{!}{\{\mappDopplOmega\mappDoppl\}}%
\resizebox{0.2in}{!}{\{\mappDopplOmega\mappDoppl\}}%
\resizebox{0.1in}{!}{\{\mappDopplOmega\mappDoppl\}}%

\hspace{0.7in}
\(=\)
\{\mappDopplOmega\}%
\resizebox{0.6in}{!}{\{\mappDoppl\mappDopplOmega\}}%
\resizebox{0.5in}{!}{\{\mappDoppl\mappDopplOmega\}}%
\resizebox{0.4in}{!}{\{\mappDoppl\mappDopplOmega\}}%
\resizebox{0.3in}{!}{\{\mappDoppl\mappDopplOmega\}}%
\resizebox{0.2in}{!}{\{\mappDoppl\mappDopplOmega\}}%
\resizebox{0.1in}{!}{\{\mappDoppl\mappDopplOmega\}}%
\(=\)
\mappDopplOmegaSquared \(=\omega^2\)
\end{itemize}

% Finally, \textbf{oridnal exponentiation} follows the formula
% \(\omega^{n+1}=\omega^n\cdot\omega\). That means \(\omega^2\) is formed by
% each Doppl in \(\omega\) splitting into its own \(\omega\),
% and \(\omega^3\) is formed
% by each Doppl in \(\omega\) splitting into \(\omega^2\)...

Somehow you're not too shocked when your \mappMobidex{} tells you that
there's a \textbf{puzzle} associated with counting Doppl. It's said
that a Doppl's favorite \mappMobimon{} attack is ``\(24007042951\)'',
but you'll need to solve the following Doppl arithmetic problems to
figure out what that means in English!

\begin{align*}
  (\omega+3)\cdot(\omega+5)&=\omega^2\cdot E+\omega\cdot G+S \\
        %M=1,G=5,S=3
  \omega+1+\omega+3+\omega+5+\omega+7&=\omega\cdot I+O \\
        %I=4,O=7
  3\cdot\omega+\omega^2\cdot 5+4\cdot(\omega^2+2)&=
        \omega^2\cdot L+\omega\cdot R+N \\
        %L=6,R=0,N=8
  2\cdot(2+\omega\cdot3)+(\omega\cdot3+2)\cdot2&=
        \omega\cdot A+M
        %A=9,M=2
\end{align*}

%solution: MIRRORIMAGE

% Include below for aucTeX integration
%%% Local Variables:
%%% mode: latex
%%% TeX-master: "../mapp-challenge-18-game-book"
%%% End:
